\section{Introduction}\label{sec:introduction}

A central charge assigns a complex number to each object of a derived
category. A stability condition organizes the objects by its argument and
requires every object to have a finite filtration by semistable factors.
Bridgeland introduced this structure and its deformation theory
\cite{Bridgeland2007}. On a projective threefold, the expected large-volume
charges are exponential expressions in the Chern character. The
prescribed-charge problem asks whether these particular expressions admit
a compatible heart, finite semistable filtrations, and a support bound
controlling the numerical classes of all semistable objects.

We construct stability conditions with the exact ordinary and
square-root-Todd large-volume charges on every smooth projective complex
threefold with trivial canonical bundle. One volume threshold works on an
open set of real twists and real ample directions. The support property
holds on the full numerical Grothendieck group, and the ordinary charge
has the standard double-tilt heart. An independent argument proves strong
tilt inequalities on every smooth projective complex threefold: a fixed
rational correction works at every volume, and the correction vanishes
beyond a threshold independent of the object and of the twist along a
fixed polarization.

All varieties are over $\CC$, and all Chern characters and intersection
products are numerical. Write $N^1(X)_{\RR}$ for the space of numerical
real divisor classes, $\operatorname{Amp}(X)$ for its ample cone, and
$N_1(X)_{\QQ}$ for the space of numerical rational curve classes.
A threefold is smooth, connected and projective.
The assertions for a disjoint union follow by applying the results to
its finitely many connected components and taking the largest constants.
Write $\Db(X)=\Db\Coh(X)$ and
$\ch^B(E)=e^{-B}\ch(E)$ for a real numerical divisor class $B$.

\subsection{Prescribed charges and the full numerical group}

The Euler pairing is
$\chi(E,F)=\sum_j(-1)^j\dim\Ext^j(E,F)$, and we use
\[
 \Knum(X)=K_0(X)/\{v:\chi(v,w)=0\text{ for every }w\in K_0(X)\}.
\]
A numerical stability condition consists of a charge homomorphism
$Z:\Knum(X)\to\CC$ and a locally finite slicing $\mathcal P$ of
$\Db(X)$. A slicing gives the semistable objects of each real phase,
with $Z(E)=m(E)e^{i\pi\phi}$, $m(E)>0$, on $\mathcal P(\phi)$,
the shift and Hom-vanishing rules, and finite Harder--Narasimhan
filtrations. Local finiteness requires the categories of sufficiently
short phase intervals to have finite length. We require the support
property of Kontsevich--Soibelman
\cite[Sections~1.2 and~2.1]{KontsevichSoibelman2008} on the full real space
$\Knum(X)_{\RR}$: for a norm there is $C>0$ such that
$\|[E]\|\leq C|Z(E)|$ for every semistable $E$.
Its heart is $\mathcal P(0,1]$. We call a stability condition geometric
if all point sheaves $\OO_x$ are stable of phase $1$.

We specify the tilt conventions needed to state the heart in our theorem.
For a real ample class $H$, a real divisor class $B$, and a coherent sheaf
of positive rank put
\[
 \mu_{H,B}(F)=\frac{H^2\ch_1^B(F)}{H^3\ch_0(F)},
\]
and give a nonzero torsion sheaf slope $+\infty$.
The slope Harder--Narasimhan filtration has semistable factors of strictly
decreasing slopes; write $\mu^+$ and $\mu^-$ for its largest and smallest
slopes. The slope torsion pair and its tilt are
\[
 \mathcal T_{H,B}=\{F:\mu^-_{H,B}(F)>0\},\qquad
 \mathcal F_{H,B}=\{F:\mu^+_{H,B}(F)\leq0\},\qquad
 \mathcal B_{H,B}=\langle\mathcal F_{H,B}[1],\mathcal T_{H,B}\rangle,
\]
with the zero sheaf in both subcategories. Here a heart is the abelian
category specified by a bounded $t$-structure, $[1]$ is the cohomological
shift, and brackets denote extension closure. Thus an object of
$\mathcal B_{H,B}$ has ordinary cohomology only in degrees $-1,0$, in the
indicated two subcategories.

For real $B$, real ample $H$ and $t>0$, the ordinary physical charge is
\begin{equation}\label{eq:intro-physical}
 Z_{B,tH}(E)=
 -\ch_3^B(E)+\frac{t^2}{2}H^2\ch_1^B(E)
 +i\left(tH\ch_2^B(E)-\frac{t^3}{6}H^3\ch_0(E)\right).
\end{equation}
The first tilt is unchanged if the polarization $H$ is replaced by
$tH$. On this first tilt, use the second slope
\[
 \nu_{B,tH}(E)=
 \frac{tH\ch_2^B(E)-t^3H^3\ch_0(E)/6}{t^2H^2\ch_1^B(E)}.
\]
Its denominator-zero value is $+\infty$. An object is semistable for
this slope if every nonzero proper subobject in $\mathcal B_{H,B}$ has
slope at most that of the quotient; stability uses a strict inequality.
The Harder--Narasimhan property holds for these real ample parameters
\cite[Section~2]{BMS2016}. Let $\mathcal T'_{B,tH}$
have strictly positive smallest Harder--Narasimhan slope, and
$\mathcal F'_{B,tH}$ have nonpositive largest slope. The standard
double-tilt heart is
\[
 \mathcal A_{B,tH}=
 \langle\mathcal F'_{B,tH}[1],\mathcal T'_{B,tH}\rangle.
\]
These definitions use the same strict-positive/nonpositive boundary
at both tilts.

\begin{theorem}[Exact large-volume charges on a real sector]
\label{thm:geometric-sector}
Let $X$ be a smooth projective complex threefold with
$K_X\simeq\OO_X$. For any $B_*\in N^1(X)_{\RR}$ and
$H_*\in\operatorname{Amp}(X)$ there are open neighborhoods $V$ of
$B_*$ and $W$ of $H_*$ and one $T>0$ such that, for all
$B\in V$, $H\in W$ and $t\geq T$, the following hold.
\begin{enumerate}[label=\textup{(\roman*)}]
\item $(Z_{B,tH},\mathcal A_{B,tH})$ is a geometric numerical
stability condition with support on $\Knum(X)_{\RR}$.
\item There is a geometric numerical stability condition, with support
on the same full numerical space, whose charge is exactly
\[
 Z^{\mathrm{LV}}_{B,tH}(E)
 =-\int_Xe^{-B-itH}\ch(E)\sqrt{\td(X)},
 \qquad \sqrt{\td(X)}=1+\frac{c_2(X)}{24}.
\]
Its heart is the one constructed in Section~\ref{sec:ambient}.
\end{enumerate}
\end{theorem}

Only triviality of $K_X$ is assumed; the additional vanishings
$H^1(X,\OO_X)=H^2(X,\OO_X)=0$ in the strict Calabi--Yau convention
are unnecessary. The quantifiers give one threshold for both charges
throughout the same real sector. The two hearts are specified separately:
the ordinary charge uses the double tilt defined above, while the
Todd-corrected charge uses the heart constructed in
Section~\ref{sec:ambient}.

The ordinary and square-root-Todd charges have large-volume antecedents in
Bayer's polynomial stability conditions
\cite[Theorem~3.2.2 and Section~4]{BayerPolynomial2009}. Polynomial stability
orders phases asymptotically as the volume tends to infinity. Our theorem
constructs a numerical Bridgeland stability condition at each finite
volume in the stated sector. The passage to finite volume must also
control the charge kernel: on a threefold of Picard rank greater than
one, the full numerical group contains classes invisible to the four
intersection degrees along a single polarization.

\subsection{Independent uniform tilt inequalities}

The numerical problem concerns the third Chern character, which is absent
from the classical Bogomolov--Gieseker inequality. Bayer--Macr\`i--Toda
introduced the double-tilt construction and proposed a third-character
inequality that would make it a source of stability
conditions~\cite{BMT2014}. Bayer--Macr\`i--Stellari established that
inequality for abelian threefolds and Calabi--Yau threefolds of abelian
type, and developed the discriminants, deformation and wall framework used
below~\cite{BMS2016}. For this part fix an ample integral divisor $H$
and a rational numerical divisor $B_0$.

For $a>0$, $b\in\RR$ and $B=B_0+bH$, the tilt slope is
\begin{equation}\label{eq:intro-tilt}
 \nu_{a,b}(E)=
 \frac{H\ch_2^B(E)-\tfrac12a^2H^3\ch_0(E)}
      {H^2\ch_1^B(E)},
\end{equation}
with value $+\infty$ when the denominator is zero.
We use the same subobject--quotient comparison for tilt stability.
At finite slope, semistability is equivalently the comparison with the
slope of the object itself. The treatment of zero-dimensional quotients
and the finite factor filtrations needed below is given in
Section~\ref{sec:tilt}.

\begin{theorem}[Uniform strong inequality]\label{thm:uniform-inequality}
Let $X$ be a smooth projective complex threefold, $H$ an ample integral
divisor, and $B_0\in N^1(X)_{\QQ}$.
There are constants $k\in\QQ_{\geq0}$ and $R_*>0$, depending only on
$(X,H,B_0)$, such that every finite-slope $\nu_{a,b}$-semistable
$E\in\mathcal B_{H,B_0+bH}$ with $\nu_{a,b}(E)=0$ satisfies
\begin{equation}\label{eq:intro-corrected}
 \ch_3^{B_0+bH}(E)
 \leq \left(\frac{a^2}{6}+k\right)H^2\ch_1^{B_0+bH}(E)
 \qquad(a>0,\ b\in\RR).
\end{equation}
For the same fixed data, the stronger bound
\begin{equation}\label{eq:intro-tail}
 \ch_3^{B_0+bH}(E)
 \leq \frac{a^2}{6}H^2\ch_1^{B_0+bH}(E)
\end{equation}
holds whenever $a>R_*$, for every $b\in\RR$.
\end{theorem}

There is no canonical-bundle hypothesis in this theorem. Both constants
are independent of the rank and Chern character of $E$. The curve class
$\Gamma=kH^2\in N_1(X)_{\QQ}$ gives the correction
$\Gamma\cdot\ch_1^{B_0+bH}(E)$ and satisfies $\Gamma\cdot H\geq0$.

Bernardara--Macr\`i--Schmidt--Zhao proved a corrected inequality of this
kind for Fano threefolds polarized by a divisor proportional to $-K_X$,
and asked for such a curve-class correction
on every polarized threefold
\cite[Theorem~1.1 and Question~2.4]{BMSZ2017}.
Martinez--Schmidt stated the version allowing a fixed rational
transverse twist \cite[Question~2.6]{MartinezSchmidt2019}.
Theorem~\ref{thm:uniform-inequality} answers these corrected-inequality
questions affirmatively and also gives a volume threshold beyond which
the correction can be removed.

The uncorrected inequality at every parameter is false. Schmidt's
blow-up example~\cite[Theorem~3.1]{Schmidt2017} and the contracted-divisor
and Weierstrass examples of Martinez--Schmidt
\cite[Theorem~1.1]{MartinezSchmidt2019} give explicit violations.
Section~\ref{sec:counterexamples} computes their bounded volume ranges
and proves a uniform bound for semistable sheaves of arbitrary rank on
the reduced exceptional surface. These comparisons retain both parts of
Theorem~\ref{thm:uniform-inequality}: the corrected all-volume statement
and the uncorrected tail.

The physical volume and the tilt variable satisfy $t=\sqrt3a$.
Consequently the strong bound \eqref{eq:intro-tail} has coefficient
$t^2/18$, whereas positivity of the ordinary charge
\eqref{eq:intro-physical} on zero-slope objects supplies the weaker
coefficient $t^2/2$. The proof of the strong inequality is independent of
the construction of the charges and their full numerical support.

The uncorrected tail also gives an all-slope quadratic formulation.
For $v=(v_0,v_1,v_2,v_3)$ define
\begin{equation}\label{eq:intro-quadratic}
 Q_{\alpha,\beta}(v)=
 \alpha(v_1^2-2v_0v_2)
 +\beta(3v_0v_3-v_1v_2)+2v_2^2-3v_1v_3.
\end{equation}
\begin{corollary}[A parabolic region with trivial correction]
\label{cor:quadratic-tail}
For every polarized smooth projective threefold $(X,H)$ there is $R_*>0$
such that, if
\[
 \alpha>f(\beta):=\frac{\beta^2+R_*^2}{2},
\]
then $Q_{\alpha,\beta}(v_H(E))\geq0$ for every nonzero semistable object
in $\mathcal B_{H,\beta H}$ for the slope
\[
 \frac{v_2-\beta v_1+(\beta^2-\alpha)v_0}{v_1-\beta v_0},
 \qquad
 v_H(E)=(H^3\ch_0(E),H^2\ch_1(E),H\ch_2(E),\ch_3(E)),
\]
again interpreted as $+\infty$ at zero denominator.
\end{corollary}
The function $f$ is continuous. This is the generalized quadratic
formulation with trivial correction: $v_H$ uses the ordinary,
unmodified Chern character. The slope in the corollary
is exactly~\eqref{eq:intro-tilt} with
$a=\sqrt{2\alpha-\beta^2}$ and $b=\beta$; Section~\ref{sec:wall}
proves the all-slope implication, including zero denominator.

\subsection{Previous constructions}

General existence of stability conditions on projective varieties is
supplied by Li's construction~\cite{Li2026} and the framework of
Li--Liu--Liu--Macr\`i--Perry--Stellari--Zhao~\cite{LiEtAl2026}.
Cheng establishes full numerical support and general nonemptiness
in~\cite{Cheng2026}. Thus nonemptiness alone is already known in the
setting of Theorem~\ref{thm:geometric-sector}. Our construction uses the
product, restriction and deformation methods of these works to prescribe
the two charges and the real sector.
The product-embedding route to full numerical support in
Section~\ref{sec:ambient} adapts the construction of
Giovenzana--Robotis--Rota--Zuliani~\cite{GiovenzanaRobotisRotaZuliani2026},
distinct from the quantitative product inputs cited there.

There is also a preceding large-volume construction with a different
character. Cheng--Feyzbakhsh identify a large-volume double-tilt heart for
the character $\td(N_{X/\mathbb P^n})^{-1}\ch$, where $N_{X/\mathbb P^n}$
is the normal bundle of their projective embedding
\cite[Section~3]{ChengFeyzbakhsh2026}. A reparameterization expresses
that heart as an ordinary double tilt while leaving a curve-class
correction in its charge.
The real-parameter extension discussed in their Section~3.3 follows from
\cite[Theorem~4.1, Proposition~4.5 and Theorem~6.2]{LiEtAl2026}.
Combining the mass--Hom estimate
\cite[Theorem~7.5]{LiEtAl2026} with Cheng's full-support criterion
\cite[Theorem~2.1]{Cheng2026} gives full numerical support for this
fixed-polarization real family, pointwise in its parameters.
Theorem~\ref{thm:geometric-sector} prescribes the ordinary and
square-root-Todd charges on an open set of real twists and ample
directions with one volume threshold. Its heart comparison adapts the
surface-section method of Cheng--Feyzbakhsh to the ordinary character
and positive weighted sections.

Restriction to lower-dimensional varieties has also led to strong tilt
inequalities. Feyzbakhsh--Koseki--Liu--Rekuski derive a corrected tilt
inequality from an improved sheaf Bogomolov--Gieseker bound, and obtain
that bound on families of Calabi--Yau threefolds from Brill--Noether
estimates on curves inside surface sections
\cite[Theorems~1.1, 1.3 and~1.4]{FeyzbakhshEtAl2025}.
Our numerical argument likewise uses successive restrictions, but its
input is a rank-uniform cohomology estimate on surfaces. A Frobenius
limit then gives the inequality at one fixed positive volume on an
arbitrary smooth projective threefold. The wall argument builds on the
discriminant-descent method of
\cite[Lemma~2.7]{BMSZ2017}; the fixed-volume defect estimate is what
makes its resulting threshold uniform in the twist and the object.

\subsection{Two independent proof strategies}

We prove the prescribed-charge theorem first, in
Sections~\ref{sec:ambient}--\ref{sec:heart}, and then establish the
strong inequalities in
Sections~\ref{sec:tilt}--\ref{sec:wall}. Figure~\ref{fig:proof} shows
the two arguments. Their distinction is mathematical: controlling four
Chern-character degrees along one polarization does not itself control
the full numerical group or provide a common threshold when the ample
direction varies.

For the charge construction, choose very ample line bundles whose
positive span contains the ample directions under consideration. They
give an embedding $i:X\hookrightarrow Y$ into a product of projective
spaces. In dimension three, products of divisor classes span all
numerical Chern-character classes; choosing the line bundles to span
$N^1(X)_{\RR}$ therefore makes $i_*$ injective on the numerical
Grothendieck group. Ambient support will then control every numerical
class on $X$.

The ambient construction must also permit a deformation which is uniform
at large volume. Liu's product-with-a-curve
construction~\cite{LiuProduct2021}, the positive-genus product theory
of~\cite{ProductsCurves2025}, the support induction
of~\cite{LiEtAl2026}, and Cheng's two-block
construction~\cite{Cheng2026} supply its starting framework. We prove
a weighted support estimate for finitely many projective blocks after
rescaling each Chern-character coordinate by its power of the volume.
In these coordinates the support constant is independent of volume, and
every fixed positive-degree correction to the exponential charge tends
to zero as the volume grows. This permits the ambient deformation
throughout an open set of twists and ample directions. A fixed polynomial
correction cancels the Todd factors in Grothendieck--Riemann--Roch, so
restriction gives exactly the chosen character on $X$.

This construction already gives full numerical support and stable point
sheaves. For the ordinary charge, compatible restrictions to smooth
surfaces identify its heart. The induced surface hearts control the
outer cohomology sheaves; positive weights combine their slope bounds
for a real ample direction. A separate argument at the real axis fixes
the zero-slope boundary of the second tilt. The Todd-charge construction
uses the geometric heart obtained by restriction.

The independent numerical proof is concentrated at one fixed volume $A$.
For a slope-zero object $E$ at this volume, put
$d=H^2\ch_1^B(E)/H^3$ and $r=\ch_0(E)$. The tilt discriminant
inequality gives $d\geq A|r|$. We prove
\[
 \frac{\ch_3^B(E)}{H^3}-\frac{A^2d}{6}\le C(d-A|r|)
\]
with $C$ independent of $E$ and $b$. The error $e=d-A|r|$ can vanish
even at large rank, so the estimates must remain proportional to this
error without an additional rank term. Truncating the ordinary slope
filtrations separates sheaf factors in a narrow slope interval from a
residual complex whose rank and first two intersection degrees are
bounded by $e$.
Langer's restriction inequality controls the resulting slope variance
\cite{Langer2004,Langer2004Addendum}. Hermitian--Einstein theory and the
Bochner identity~\cite{Donaldson1985,UhlenbeckYau1986,DemaillyCADG},
together with a matrix spectral estimate of Cwikel--Lieb--Rozenblum type
\cite{Frank2014}, control the residual cohomology in terms of $e$.
For each fixed object, we reduce finitely many
sheaves and maps to positive characteristic. Frobenius pullback and
Riemann--Roch turn the cohomology bounds into the displayed inequality
on $X$; the errors depending on the chosen object vanish in the limit.
Transport along zero-slope branches, with strict discriminant descent at
walls, gives the corrected global bound and the uncorrected tail.

\begin{figure}[htbp]
\centering
\begin{tikzpicture}[
 box/.style={draw,rounded corners=2pt,align=center,font=\small,
             text width=4.4cm,minimum height=10mm,inner sep=5pt},
 arrow/.style={-{Stealth[length=2mm]},thick}]
\node[box] (ambient) {Weighted ambient stability\\and support};
\node[box,below=7mm of ambient] (deform) {Exact charge deformation\\and restriction to $X$};
\node[box,below=7mm of deform] (support) {Full numerical support\\and stable points for both charges};
\node[box,below=7mm of support] (heart) {Ordinary charge: surface comparison\\and the double-tilt heart};
\node[box,right=15mm of ambient] (analysis) {Rank-uniform restriction\\and cohomology};
\node[above=3mm of ambient,align=center,font=\small] {Prescribed charges: $K_X\simeq\OO_X$};
\node[above=3mm of analysis,align=center,font=\small] {Inequalities: arbitrary threefold};
\node[box,below=7mm of analysis] (apex) {Fixed-volume defect bound};
\node[box,below=7mm of apex] (wall) {Wall transport};
\node[box,below=7mm of wall] (bounds) {Corrected all-volume inequality\\and uncorrected tail};
\draw[arrow] (ambient)--(deform);
\draw[arrow] (deform)--(support);
\draw[arrow] (support)--(heart);
\draw[arrow] (analysis)--(apex);
\draw[arrow] (apex)--(wall);
\draw[arrow] (wall)--(bounds);
\end{tikzpicture}
\caption{The independent categorical and numerical arguments. Full support
is obtained in Section~\ref{sec:ambient}, before the ordinary-heart
comparison in Section~\ref{sec:heart}. The numerical argument begins
with the tilt tools of Section~\ref{sec:tilt} and proves the inequalities
in Sections~\ref{sec:analysis}--\ref{sec:wall}. The relation $t=\sqrt3a$
compares volume conventions.}
\label{fig:proof}
\end{figure}

For the numerical argument, Section~\ref{sec:tilt} supplies the circle
and duality tools, Sections~\ref{sec:analysis} and~\ref{sec:apex}
establish the fixed-volume bound, and Section~\ref{sec:wall} proves
Theorem~\ref{thm:uniform-inequality} and
Corollary~\ref{cor:quadratic-tail}.
Section~\ref{sec:counterexamples} compares the inequalities with
counterexamples and bounds arbitrary-rank sheaves on reduced exceptional
surfaces.
